\section{Discussion}

The experiments identify three decisions that shape multi-teacher distillation: which responses receive loss weight, how their gradients combine with optimizer history, and which increments survive model writeback. At a common early checkpoint, the averaging rule changes the capability tradeoff. At a fixed student state, input composition changes teacher-conditioned directions, and broader vocabulary supervision changes gradient direction even when BF16 activity remains nearly constant. Reporting these quantities together makes parameter geometry more informative about the optimization process.

The present evidence covers one model family, four related teachers, one training seed, and local calculations at update 100. The two diagnostic banks contain eight responses generated with a 256-token cap; all reach that cap, and the selected prefixes have high top-64 coverage. Extending the analysis to later and higher-entropy prefixes would test how this coverage affects the supervision comparison. Completing the aligned checkpoint grid and adding training seeds would establish the persistence of the capability tradeoffs. Restricted-parameter training would connect the measured gradient supports to operational subnetworks.
